\documentclass[11pt,a4paper]{article}
\usepackage[utf8]{inputenc}
\usepackage[T1]{fontenc}
\usepackage[ngerman]{babel}
\usepackage[margin=2.4cm]{geometry}
\usepackage{booktabs}
\usepackage{longtable}
\usepackage{xcolor}
\usepackage[colorlinks=true,urlcolor=blue!60!black,linkcolor=black]{hyperref}
\usepackage{parskip}
\usepackage{enumitem}
\setlist{nosep,leftmargin=*}

\newcommand{\code}[1]{\texttt{\small #1}}
\newcommand{\ok}{\textcolor{green!45!black}{\textbf{erledigt}}}
\newcommand{\offen}{\textcolor{orange!75!black}{\textbf{offen}}}
\newcommand{\upstream}{\textcolor{blue!60!black}{\textbf{upstream}}}

\title{AWI-ESM3-4-2-veg-HR, CMIP7 piControl\\[2mm]
\large Antwort auf deinen Koordinaten-Check von \code{cli112}}
\author{Jan Streffing, AWI}
\date{13. August 2026}

\begin{document}
% babel-ngerman macht " aktiv, was in Code-Beispielen zuschlaegt:
% operation="instant" wurde als operation=ïnstant" gesetzt. Die Shorthands
% werden erst bei \begin{document} scharfgeschaltet, also muss das hierhin.
\shorthandoff{"}
\maketitle

\section{Kurzfassung}

Danke f\"ur den Koordinaten-Check. Wir haben alle deine Punkte auf den 540 Dateien
von \code{cli112} nachgestellt, jeder einzelne reproduziert. Sieben von acht sind
behoben, einer ist offen und hat eine Frage an dich dran.

Ein Unterschied zum letzten Mal: die Zahlen unten sind der Befund auf
\code{cli112}, nicht das Ergebnis danach. Der Lauf, der die Korrekturen zeigt,
l\"auft gerade. Wir schicken die Gegenprobe nach, statt sie hier schon zu
behaupten.

\begin{center}
\begin{tabular}{llr l}
\toprule
& \textbf{Befund} & \textbf{Dateien} & \textbf{Stand} \\
\midrule
1 & unstrukturiert, ganz ohne \code{lat}/\code{lon}      & 16  & \ok \\
2 & \code{lev}-Dimension ohne \code{lev}-Variable        & 2   & \ok \\
3 & horizontal vor vertikal                              & 25  & \ok \\
4 & Dimension hei\ss t \code{plev19} / \code{plev3}      & 17  & \ok \\
5 & \code{lat\_bnds} \"uber \code{nvertex}               & 447 & \ok \\
6 & \code{coordinates} listet Koordinatenvariablen       & 490 & \ok \\
7 & \code{typetree*} alle mit Wert \code{trees}          & 4 Regeln & \ok \\
8 & \code{tclm} ohne \code{climatology\_bnds}            & 1   & \offen \\
\bottomrule
\end{tabular}
\end{center}

\section{Die einzelnen Punkte}

\subsection*{1. Unstrukturierte Dateien ganz ohne \code{lat}/\code{lon}}
\ok\ 16 Dateien hatten nur einen nackten Zellindex und keine geographische
Koordinate: \code{rlds}, \code{rlus}, \code{rsds}, \code{rsus}, \code{sbl},
\code{sifllattop}, \code{siflsenstop}, \code{volcello} und \code{deptho} auf dem
Knotengitter, \code{hfx} und \code{hfy} auf dem Elementgitter.

Die Ursache ist unspektakul\"ar und daf\"ur lehrreich: die meisten FESOM-Streams
bringen \code{lat}/\code{lon} von XIOS schon mit, deshalb steht in keinem unserer
Rezepte ein expliziter Gitter-Schritt. Bei den paar Streams, die sie nicht
mitbringen, f\"allt das entsprechend still hinten runter.

Wir holen sie jetzt aus dem Mesh, auf das ohnehin jede FESOM-Regel zeigt. Knoten
direkt, Elemente aus der \code{triag\_nodes}-Konnektivit\"at: die drei Eckknoten
sind die Zellecken, ihr Mittel der Schwerpunkt. Gemittelt wird \"uber
Einheitsvektoren, sonst landet ein Dreieck auf der Datumsgrenze bei
L\"ange 0 statt 180. Gegengepr\"uft: alle 6229020 Schwerpunkte liegen im
Breitenbereich ihres eigenen Dreiecks.

\subsection*{2. \code{lev}-Dimension ohne \code{lev}-Variable}
\ok\ Betroffen waren \code{masscello} und \code{thkcello}, und da steckte mehr
dahinter als eine fehlende Koordinate. Beide kamen als reines Profil
\code{(lev)} heraus, ohne jede horizontale Dimension, obwohl ihre Branding
\code{ti-ol-hxy-sea} ein horizontales Feld verlangt. Die zeitabh\"angigen
Geschwister derselben Variablen haben \code{(time, nod2, lev)}.

Beim Nachsehen ist \code{volcello} als \code{fx} gleich mit aufgefallen, das du
nicht gemeldet hattest: es lag als \code{volcello(ncells, nz1)} vor.
\code{nz1} ist ein FESOM-interner Name, eine Level-Koordinate gab es auch dort
nicht. Alle drei liegen jetzt als \code{(lev, ncells)} vor.

Zwei Sachen sind dabei mit herausgefallen, die wir sonst nicht gefunden h\"atten:

\begin{itemize}
\item Unser Mesh deklariert 57 Level, aber die letzten Eintr\"age sind Datenm\"ull.
      \code{depth} endet auf 5825, 6125, \emph{3160}, \code{depth\_bnds} auf 6000,
      6250, \emph{70}. Echt sind 56 Schichten, was FESOM selbst auch so schreibt.
      Deshalb waren die \code{fx}-Dateien eine Ebene l\"anger als die
      Monatsdateien, die sie beschreiben sollen. Wir schneiden jetzt auf den
      monoton steigenden Anfang zu, statt der deklarierten L\"ange zu glauben.
\item Die S\"aulen endeten nicht am Meeresboden. Das tun sie jetzt, und die
      Abbruchbedingung ist gemessen statt geraten: ein Knoten mit
      \code{depth\_lev = 3} hat im Modellstrom 4 nasse Ebenen, also ist Ebene $k$
      nass f\"ur $k \le$ \code{depth\_lev}. Stimmt bei allen Stichproben, nasser
      Anteil 64{,}1\,\%, Dicken 5 bis 350\,m wie im Modellstrom.
\end{itemize}

\subsection*{3. Horizontal vor vertikal}
\ok\ Hier hattest du recht und wir hatten den Block zuerst falsch eingeordnet.
Wir hatten ihn f\"ur ein Artefakt des Checkers gehalten, weil er die
unstrukturierte Zelldimension als horizontale Achse vom Typ A r\"at. Das stimmt
auch f\"ur einen Teil, aber eben nur f\"ur einen: 25 Dateien haben die
Reihenfolge wirklich verkehrt herum, \code{difvso(time, nod2, lev)},
\code{hfx(time, elem, lev)} und so weiter. FESOM schreibt die Vertikale hinten,
und daran hat nach dem Dimension-Mapping niemand mehr gedreht.

Wir sortieren die Vertikale jetzt vor die erste horizontale Dimension. Bewusst
nur das: CF stellt auch die \"ubrigen Dimensionen gern nach links, aber die
Dateien laufen dort durch, und wir wollten nicht Ausgabe umr\"uhren, zu der es
keinen Befund gibt.

\subsection*{4. \code{plev19} und \code{plev3} als Dimensionsname}
\ok\ In \code{CMIP7\_coordinate.json} haben \code{plev3}, \code{plev19} und
\code{plev39} alle den \code{out\_name} \code{plev}. Die Zahl ist eine Stufe des
Data Request und geh\"ort in die L\"ange der Dimension, nicht in ihren Namen.
Sonst \"andert sich der Dimensionsname mit der angeforderten Stufe. Betraf 17
Dateien \"uber \code{hur}, \code{hus}, \code{ta}, \code{ua} und \code{va}.

\subsection*{5. Zellgrenzen: \code{vertices\_latitude} statt \code{lat\_bnds}}
\ok\ Und der Befund war schlimmer, als er in deiner Liste aussieht, weil unsere
Ausgabe nicht einmal in sich konsistent war. Je nachdem, ob die Grenzen aus dem
XIOS-Strom oder aus dem Mesh kamen, hie\ss\ die Vertex-Dimension anders:

\begin{center}
\begin{tabular}{lr}
\toprule
\textbf{Vertex-Dimension} & \textbf{Dateien} \\
\midrule
\code{nvertex}  & 330 \\
\code{vertices} & 117 \\
\bottomrule
\end{tabular}
\end{center}

Jetzt einheitlich \code{vertices\_latitude(<Zellen>, vertices)} und
\code{vertices\_longitude}, wie CMOR. Nur f\"ur echte Hilfskoordinaten, also
unstrukturierte Gitter; auf regul\"aren Gittern bleibt \code{lat\_bnds(lat, bnds)}
stehen, das ist dort richtig.

\subsection*{6. \code{coordinates} listet Koordinatenvariablen}
\ok\ Wir hatten das Attribut aus allen Dimensionen und Koordinaten der Variablen
zusammengebaut, statt nur aus den Hilfskoordinaten. CF \S5 will dort das, was
\"uber die Dimensionsnamen nicht auffindbar ist. So sah das aus:

\begin{center}
\begin{tabular}{rl}
\toprule
\textbf{Dateien} & \textbf{\code{coordinates}} \\
\midrule
199 & \code{time lat lon} \\
115 & \code{time lon lat} \\
 40 & \code{time} \\
 30 & \code{time lev lat lon} \\
 27 & \code{time lon lat type} \\
 23 & \code{time lat lon height} \\
 19 & \code{time lon lat sector} \\
\bottomrule
\end{tabular}
\end{center}

\code{lat} und \code{lon} geh\"oren auf dem unstrukturierten Gitter da hin, auf
dem regul\"aren nicht, und \code{time} und \code{lev} in keinem Fall.
\code{type}, \code{sector} und \code{height} sind Skalare und bleiben. Wo nach
dem Aufr\"aumen nichts \"ubrig bleibt, f\"allt das Attribut ganz weg.

\subsection*{7. \code{typetree*}}
\ok\ Und hier w\"are ein Blick von deiner Seite n\"utzlich. Unsere Kopie von
\code{CMIP7\_coordinate.json} tr\"agt bei allen vier Bl\"atter-Varianten den Wert
\code{trees}:

\begin{center}
\begin{tabular}{lll}
\toprule
\textbf{Variable} & \textbf{geschrieben} & \textbf{richtig w\"are} \\
\midrule
\code{treeFrac}       & \code{trees} & \code{trees} \\
\code{treeFracBdlDcd} & \code{trees} & \code{broadleaf\_deciduous\_trees} \\
\code{treeFracBdlEvg} & \code{trees} & \code{broadleaf\_evergreen\_trees} \\
\code{treeFracNdlDcd} & \code{trees} & \code{needleleaf\_deciduous\_trees} \\
\code{treeFracNdlEvg} & \code{trees} & \code{needleleaf\_evergreen\_trees} \\
\bottomrule
\end{tabular}
\end{center}

Damit sind die vier Variablen auf der Platte nicht auseinanderzuhalten, obwohl
die Blattform der ganze Punkt an ihnen ist. Die CF-Area-Type-Tabelle f\"uhrt die
spezifischen Begriffe seit langem, und alle anderen Eintr\"age in derselben
Tabelle sind bei uns korrekt spezifisch, etwa \code{cropFracC3} mit
\code{crops\_of\_c3\_plant\_functional\_types}.

CMIP6 hatte an derselben Stelle ebenfalls \"uberall \code{trees}. Unsere Kopie
tr\"agt \code{table\_date 2026-07-09}. Wir setzen die vier Werte lokal, mit
einer Notiz, das wieder rauszunehmen, sobald es upstream steht. Falls die
ver\"offentlichte Tabelle das noch so f\"uhrt, w\"are das vermutlich f\"ur mehr
Gruppen als uns relevant.

\subsection*{8. \code{tclm} ohne \code{climatology\_bnds}}
\offen\ Du hast recht, die Datei schreibt \code{time\_bnds}. Wir m\"ochten das
aber nicht einfach umbenennen, weil dann etwas Falsches sauber aussieht.

Betroffen ist genau eine Regel, \code{atmos.pfull.tclm-al-hxy-u.mon.glb}. Die
Pipeline dahinter bildet gar keine Klimatologie, sie mittelt ein Jahr monatsweise
und schreibt das unter dem \code{tclm}-Label. Nur die CF-Attribute
nachzur\"usten, w\"urde eine mislabelte Datei durch den Checker bringen. Richtig
ist, \"uber die Jahre zu mitteln, und das steht bei uns f\"ur den
Produktionslauf an.

Die Frage an dich: ist \code{tclm} bei monatlicher Frequenz so gemeint, dass die
Klimatologie \"uber den gesamten publizierten Zeitraum geht? Dann h\"angt die
Datei am Ende des Laufs und nicht an einem einzelnen Jahr, und der
Dateiname-Zeitraum m\"usste den ganzen Bereich abdecken.

\section{Noch offen von letztem Mal}

Danke f\"ur die Einordnung zu \code{grid\_label} bei Dateien ohne horizontales
Gitter. Wir registrieren f\"ur die zonalen und Becken-Variablen ein eigenes
Output-Grid, statt sie unter dem Gitter mitlaufen zu lassen, aus dem sie
gerechnet wurden.

Zur \code{vegtype}-Achse: sie deckt jetzt alle 44 PFT-Spalten ab, verteilt auf
elf Kategorien aus der CF-Area-Type-Tabelle. Moose und Flechten laufen unter
\code{vegetation} mit, weil es f\"ur sie keinen eigenen CF-Begriff gibt. Einen
neuen zu beantragen, halten wir nicht f\"ur unsere Aufgabe, und still wegfallen
lassen wollten wir sie erst recht nicht. Nicht zuordenbare Spalten l\"osen jetzt
einen Fehler aus, statt zu verschwinden; es gibt aktuell keine mehr.

\end{document}
